\documentclass[12pt]{article}

% Idioma e codificação
\usepackage[brazil]{babel}
\usepackage[utf8]{inputenc}
\usepackage[T1]{fontenc}

% Fonte semelhante à Arial
\usepackage{helvet}
\renewcommand{\familydefault}{\sfdefault}  % Usa Helvetica (Arial) em todo o texto

% Espaçamento e margens
\usepackage{setspace}
\onehalfspacing
\usepackage{geometry}
\geometry{a4paper, left=2cm, right=2cm, top=2cm, bottom=2cm}
\usepackage{float}


% Parágrafo
\usepackage{indentfirst}
\setlength{\parindent}{1.25cm}
\setlength{\parskip}{0pt}

% Títulos personalizados
\usepackage{titlesec}
\titleformat{\section}
  {\normalfont\sffamily\bfseries\fontsize{20}{20}\selectfont} % Arial 20 em negrito
  {\thesection}{1em}{} % remove numeração se quiser
\titleformat{\subsection}
  {\normalfont\sffamily\bfseries\fontsize{14}{14}\selectfont}
  {\thesubsection}{1em}{}

% Legendas (captions) personalizadas - Arial 10
\usepackage[font={sf,footnotesize},labelfont=bf,labelsep=colon,justification=centering]{caption}

  

\begin{document}
\begin{center}
\textbf{Fundação Escola de Comércio Álvares Penteado (FECAP)}\\
\textbf{Curso de Ciência da Computação - 4º Semestre}\\[2cm]

\vspace*{\fill}
\textbf{\Large Cálculo de Intervalo de Confiança (IC)}\\[0.5cm]
\textbf{Entrega 2 – Análise Inferencial de Dados}\\[2cm]
\end{center}


\noindent\hspace*{6cm}\textbf{Caroliny Rossi Bittencourt – 24025959}\\
\hspace*{6cm}\textbf{Rafael Alves dos Santos Guimarães – 24025724}\\
\hspace*{6cm}\textbf{Duda Lucena Miguel – 24025889}\\
\hspace*{6cm}\textbf{Isadora Teixeira Santoma – 24026458}\\[2cm]

\vspace*{\fill}

\begin{center}
\vfill
\textbf{São Paulo – 2025}
\end{center}

\newpage
\onehalfspacing

\section*{Introdução}

O intervalo de confiança (IC) é uma ferramenta estatística utilizada para estimar o intervalo no qual se encontra o valor verdadeiro de uma média ou proporção populacional, com determinado nível de confiança, geralmente 95\%.  
Em termos práticos, o IC expressa a precisão de uma estimativa: quanto menor sua amplitude, maior a confiança de que o valor obtido representa bem a população.

No contexto empresarial, o uso de intervalos de confiança é essencial para apoiar a tomada de decisões. Ele permite avaliar a variabilidade de métricas como vendas, lucros ou, neste caso, o uso de cupons por loja. Assim, gestores podem inferir tendências e níveis de desempenho sem depender apenas de observações pontuais.

\section*{Metodologia}

Foi utilizada uma base de dados fornecida pela própria empresa, \textbf{PicMoney}, contendo informações sobre o uso de cupons de diferentes estabelecimentos comerciais parceiros da plataforma.
O cálculo do intervalo de confiança foi realizado a partir da média de cupons utilizados por loja, considerando um nível de confiança de 95\%, com base na distribuição normal.

O arquivo foi analisado no mesmo diretório do presente documento, utilizando a linguagem de programação \textbf{Python}, com o auxílio das bibliotecas \textbf{Pandas} e \textbf{NumPy} para manipulação e análise estatística dos dados. O Script de códigos desse cálculo também se encontra no mesmo diretório deste documento, tendo como nome "main.ipynb" disponível para eventual consulta ou reprodução dos resultados obtidos.  
Para executá-lo, é necessário ter o \textbf{Python} instalado na máquina e utilizar um ambiente de desenvolvimento, como o \textbf{VSCode} ou outra IDE de preferência.  
Após configurar o ambiente, devem ser instaladas as bibliotecas utilizadas por meio do comando:

\begin{center}
\texttt{pip install pandas numpy}
\end{center}

Em seguida, basta abrir o arquivo do código, garantir que o diretório de trabalho esteja corretamente configurado, e executar o programa para que os cálculos sejam realizados automaticamente conforme os parâmetros apresentados neste relatório.

\subsection*{Metodo de Cálculo}
Para o cálculo do intervalo de confiança da média, foram utilizados os seguintes parâmetros:

\begin{itemize}
    \item Média amostral ($\bar{x}$)
    \item Desvio padrão ($s$)
    \item Número de observações ($n$)
    \item Valor crítico de $Z$ para 95\% de confiança (1,96)
\end{itemize}

A fórmula aplicada foi:

\[
IC = [\bar{x} - Z \times \frac{s}{\sqrt{n}} ; \bar{x} + Z \times \frac{s}{\sqrt{n}}]
\]

onde:
\begin{itemize}
    \item $\bar{x}$ representa a média amostral de cupons por loja;
    \item $s$ é o desvio padrão das observações;
    \item $n$ é o número total de lojas analisadas.
\end{itemize}

Dessa forma, obtevem-se o intervalo de confiança da média de cupons por loja com base nos dados disponíveis

\section*{Desenvolvimento}

Ao analisarmos os dados reais da amostra, observamos que a maioria das lojas apresentou valores de cupons próximos da média geral, situada em aproximadamente \textbf{3025 cupons por loja}, como exposto na Tabela 1 a seguir:

\begin{table}[H]
\centering
\begin{tabular}{|l|c|}
\hline
\textbf{Nome do Estabelecimento} & \textbf{Cupons Utilizados} \\
\hline
Açaí no Ponto & 2777 \\
Burger King & 2585 \\
Café Cultura & 3116 \\
Carrefour Express & 3073 \\
Casas Bahia & 2007 \\
Churrascaria Boi Preto & 2912 \\
Clube Pinheiros & 3079 \\
Droga Raia & 4271 \\
Drogaria São Paulo & 4310 \\
Drogasil & 4338 \\
Extra & 3090 \\
Fleury & 3087 \\
Forever 21 & 3144 \\
Habib's & 2630 \\
Just Run & 3061 \\
Lavoisier & 3185 \\
Madero & 2883 \\
Magazine Luiza & 1980 \\
McDonald's & 2699 \\
Octavio Café & 3162 \\
Outback & 2837 \\
Ponto & 1946 \\
Pão de Açúcar & 3075 \\
Renner & 3161 \\
Riachuelo & 3121 \\
Ráscal & 2914 \\
Sabin & 3217 \\
Selfit & 3098 \\
Sesc Carmo & 3120 \\
Sesc Paulista & 3103 \\
Smart Fit & 3147 \\
Starbucks & 3084 \\
Subway & 2631 \\
\hline
\textbf{Média geral} & \textbf{3025,55} \\
\hline
\end{tabular}
\caption{Quantidade de cupons utilizados por loja}
\end{table}

\vspace{0.5cm}

\newpage
\onehalfspacing

O cálculo do intervalo de confiança resultou em

\[
\textbf{IC = [2842,99 ; 3208,10]}
\]

o que indica que, com 95\% de confiança, a média verdadeira de cupons por loja encontra-se dentro desse intervalo.


\section*{Conclusão}

Com base nessa análise, foi possível chegar a uma conclusão prática aplicada ao contexto empresarial:  
embora exista certa variação entre as lojas, a maioria mantém desempenho semelhante, o que demonstra uma consistência estatística no uso dos cupons.  
Alguns estabelecimentos, como as redes de farmácias (Droga Raia, Drogaria São Paulo e Drogasil), apresentaram valores significativamente acima do limite superior do intervalo, indicando maior engajamento com a estratégia de cupons.  
Por outro lado, lojas como Casas Bahia, Ponto e Magazine Luiza ficaram abaixo do limite inferior, sugerindo menor adesão à campanha.

Esses resultados reforçam a utilidade do intervalo de confiança como instrumento para a interpretação de tendências e identificação de outliers, fornecendo base sólida para que a empresa realize ajustes estratégicos de acordo com o comportamento de cada segmento.



\end{document}
